\section{A pixel mass law in native memory dynamics}
\label{sec:mcpr-native-memory}

A chosen simulator specification can realize the conditional pixel-dependent
charged mass law on the native gauge-family carrier. The construction uses
the twelve-port reset controller and the complete native current instruments.
Its design inputs are independent left and right probe durations, a rule
enabling feedback on a successful C event, and a common feedback amplitude.
The conditional mass-law architecture is target-informed; evaluating that
architecture without a charged-mass input does not establish a prospective
prediction.

The selected response architecture gives
\begin{align}
 \kappa&=\frac{\alpha_U/50}{1+\alpha_U/31+\alpha_U^2/310},&
 \chi&=\frac{\alpha_U^2/512}{1-\alpha_U/77},\\
 \zeta&=\frac{(P-\varphi)^2/(21\pi)}{1-\alpha_U/27-\alpha_U^2/135},&
 r_j&=1+\sqrt2 e^{-\chi}
       \cos\left(\frac29+\zeta+\frac{2\pi j}{3}\right),
\end{align}
where $\varphi=(1+\sqrt5)/2$ and $j=0,1,2$.
Its charged coefficients, ordered increasingly, are
\begin{equation}
 t_i=\frac{m_i}{v}
     =\frac{e^\kappa r_i^2}
      {\left(2\,6^{14}\prod_{j=0}^{2}r_j^2\right)^{1/3}},
 \qquad \frac v{E_*}=P^{-1/2}
       \exp\left(-\frac{2\pi}{4\alpha_U}\right).
\end{equation}
The register dimensions and determinant exponent in these formulas are
part of the selected architecture. The scale $E_*$ is a source coordinate;
assigning it an energy in laboratory units requires a clock calibration.

Let $t_i>0$ be these three charged coefficients and
$\lambda_i=t_i/\sum_jt_j$. For the native projectors $P_2,P_3,P_{10}$,
direct multiplication gives $\operatorname{tr}(P_iP_j)=1/5$ for distinct
indices and the spanning-vector Gram determinant
$d=2/5-2\sqrt5/25$. Consequently the matrix
$Y=uP_2+v_1P_3+wP_{10}$ has elementary eigenvalue invariants
\begin{equation}
 \operatorname{tr}Y=u+v_1+w,\qquad
 e_2(Y)=\frac45(uv_1+uw+v_1w),\qquad
 \det Y=d uv_1w.
\end{equation}
Choose its three positive weights as the roots of
\begin{equation}
 z^3-z^2+\frac54e_2(\lambda)z-\frac{e_3(\lambda)}d=0,
 \qquad w<u<v_1.
 \label{eq:memory-shape-cubic}
\end{equation}
It follows that $\operatorname{spec}Y=\{\lambda_i\}$.

The complete C-enabled memory coefficient, with $a=p(1-q)$,
$b=pq$ and $D=1-a^2+ab$, is
\begin{equation}
 Z_C=\frac{(1-a)P_2+a(1-a)P_3+abP_{10}}D-(1-p)P_2.
\end{equation}
It has trace $p$; disabled records have not been removed or normalized away.
Put $A=v_1(v_1+w)$ and
$B=v_1(1-v_1+w)+(v_1+w)w$. The inverse equations
\begin{equation}
 -Ap^2+Bp-w=0,\qquad
 q=1-\frac{v_1}{1-p(v_1+w)}
 \label{eq:memory-rate-inverse}
\end{equation}
give the exact identity $Z_C=pY$. This follows by substituting the second
equation and reducing each of the three coefficient differences by the
quadratic in the first equation.

At each of the supplied forward and measured-endpoint pixel coordinates,
rational Taylor enclosures and polynomial endpoint signs give exactly
three roots of~\eqref{eq:memory-shape-cubic}, and exactly two solutions of
\eqref{eq:memory-rate-inverse} satisfying $0<p<1/12$, $0<q<1/6$.
For the forward coordinate the first solution is approximately
$p=0.01915526454$, $q=0.05474497127$.
The native charged-lepton support Casimirs are $9/8$ and $3$, so the
complete probes implement these rates with the declared durations
\begin{equation}
 \theta_L=\frac{\arcsin\sqrt{12p}}{\sqrt{9/8}},\qquad
 \theta_R=\frac{\arcsin\sqrt{6q}}{\sqrt3}.
\end{equation}
Both probes retain all thirteen outcomes. Their product retains all
$169\times128$ current and repair labels per old memory input.
The covariant binary grouping retains $512$ labels per input and all
twelve old inputs.

Choose $g=\sum_i t_i/p$. Then $gZ_C$ has exactly the three coefficients
$t_i$. Its diagonalizing unitary acts on the gauge-trivial family factor,
commutes with every native gauge current, and preserves the canonical
kinetic form. Thus this is an operator realization of the conditional
mass law, up to the allowed family frame. It does not assert an additional
alternating-group intertwiner. The source-coordinate enclosures also give
$\sum_it_i<1$ by bounding the cube of that sum.

The coefficient is the generator of a controlled limit of the full
correlated memory. Apply $\exp(-ihgX\otimes P_{\rm new})$ on enabled
records and identity on disabled records. Let $\Psi_h$ be this complete
classical/quantum channel. Its classical transition satisfies
$T=(1-p)R+pS$, where $R$ resets to port two and $S$ is stochastic.
For every normalized initial classical/quantum state $z_0$, including
an arbitrary untouched finite reference, the reduced state after $N$
steps of size $h=\tau/N$ obeys
\begin{equation}
 \left\|\sum_{\rm memory}\Psi_{\tau/N}^{N}(z_0)
       -\mathcal U_\tau\!\left(\sum_{\rm memory}z_0\right)\right\|_1
 \le \frac{4g\tau}{N(1-p)}
       +\left(\frac8{1-p}+4\right)\frac{g^2\tau^2}{N},
 \label{eq:memory-mass-averaging}
\end{equation}
where $\mathcal U_\tau$ is generated by $gX\otimes Z_C$.
For a stationary product preparation the first term vanishes.

To prove the bound, let $\Pi z=\pi\sum z$ and $e_k=(I-\Pi)z_k$.
The sum of trace norms on memory blocks gives
$\|Te\|_1\le p\|e\|_1$ on zero-sum blocks and
$\|\Psi_h-T\|\le2gh$. Hence
$\|e_k\|_1\le2p^k+4gh/(1-p)$.
The reduced stationary step and the mean unitary step have the same first
derivative; their two Taylor remainders sum to at most $4g^2h^2$.
The reduced error due to $e_k$ is at most $2gh\|e_k\|_1$.
Summing these errors proves~\eqref{eq:memory-mass-averaging}.

This limit is in dimensionless simulator time. A finite unread channel
need not equal its limiting unitary. The feedback pulse service is an
explicit design choice; compilation into a smaller native gate alphabet
is a separate requirement. The forward pixel coordinate has no complete
physical endpoint proof, and the measured-endpoint coordinate inherits
measured-alpha calibration. No physical pole-mass identification,
renormalization map, operational energy scale, or complete quark and
neutrino spectrum follows from this finite realization.
